\section{Dimensjonering, logistikk og operasjonsmodell}
\label{sec:dimensjonering}

Beregningene under demonstrerer at en beskjeden andel av meieriflåten har tilstrekkelig hydraulisk og logistisk kapasitet til å dekke en hel regions vannbehov.

\subsection{Matematisk omløpsmodell}
Modellen bygger på empiriske driftsparametere fra norsk tankbiltransport:
\begin{itemize}[leftmargin=1.5em, itemsep=0.15em]
    \item Tankvolum ($V$): \SI{20}{\cubic\meter} (\SI{20000}{\liter}).
    \item Avstand ($d$): \SI{30}{\kilo\meter} én vei (\SI{60}{\kilo\meter} tur/retur kilde).
    \item Gjennomsnittsfart ($v$): \SI{50}{\kilo\meter\per\hour} (inkludert flomomkjøringer).
\end{itemize}

Total syklustid per rotasjon ($T_{\text{rot}}$) defineres ved:
\begin{equation}
    T_{\text{rot}} = t_{\text{fyll}} + t_{\text{transitt}} + t_{\text{loss}} + t_{\text{buffer}}
    \label{eq:syklus}
\end{equation}
hvor leddene kvantifiseres som følger:
\begin{align*}
    t_{\text{fyll}} &= \frac{V}{Q_{\text{pumpe}}} = \frac{\SI{20000}{\liter}}{\SI{1200}{\liter\per\minute}} \approx \SI{17}{\minute}, \\
    t_{\text{transitt}} &= \frac{2d}{v} = \frac{\SI{60}{\kilo\meter}}{\SI{50}{\kilo\meter\per\hour}} = \SI{1.2}{\hour} = \SI{72}{\minute}, \\
    t_{\text{loss}} &\approx \SI{90}{\minute} \quad (\text{vektet snitt mellom Linje A og B}), \\
    t_{\text{buffer}} &\approx \SI{46}{\minute} \quad (\text{slangetilkobling, hygieneinspeksjon, pause}).
\end{align*}
Dette gir $T_{\text{rot}} = 17 + 72 + 90 + 46 = \SI{225}{\minute} = \SI{3.75}{\hour}$.

I en to-skifts døgnoperasjon gjennomfører hver bil konservativt $n = 3$ omløp per døgn, som gir $V_{\text{døgn,bil}} = 3 \cdot \SI{20000}{\liter} = \SI{60000}{\liter}$. Mobiliseres $N = 15$ melketankbiler, oppnås:
\begin{equation}
    V_{\text{total}} = N \cdot V_{\text{døgn,bil}} = 15 \cdot \SI{60000}{\liter} = \SI{900000}{\liter\per\text{døgn}}.
    \label{eq:totalvolum}
\end{equation}

Ved \textsc{dsb}s basale krisebehov på $q_{\text{basal}} = \SI{3}{\liter}$ per person per døgn dekker systemet:
\begin{equation*}
    P_{\text{dekket}} = \frac{\SI{900000}{\liter}}{\SI{3}{\liter\per\text{person}}} = \num{300000}\text{ personer}.
\end{equation*}
For en by med \num{75000} innbyggere gir dette en sikkerhetsfaktor på 4.0, eller \SI{12}{\liter} per person per døgn inkludert basal hygiene og helsestell.

\begin{table}[htbp]
\centering
\small
\caption{Dimensjonerende nøkkelparametre for Prosjekt Meierikjeden.}
\label{tab:parametre}
\begin{tabularx}{\textwidth}{lXr}
\toprule
\textbf{Parameter} & \textbf{Fysisk betydning / Grunnlag} & \textbf{Verdi} \\
\midrule
Tankvolum per bil ($V$) & Syrefast termosbeholder (\textsc{aisi} 316L) & \SI{20}{\cubic\meter} (\SI{20000}{\liter}) \\
Flåtestørrelse ($N$) & Mobiliserte biler (turnusreserver og buffer) & 15 enheter \\
Pumpekapasitet ($Q_{\text{pumpe}}$) & Sanitær impellerpumpe & \SI{1200}{\liter\per\minute} \\
Omløpstid ($T_{\text{rot}}$) & Inkluderer fylling, \SI{60}{\kilo\meter} transitt og lossing & \SI{3.75}{\hour} \\
Rotasjoner per bil ($n$) & Konservativ skiftdrift & 3 omløp/døgn \\
\textbf{Total døgnleveranse} & \textbf{Flåtens samlede sterile vannvolum} & \textbf{\SI{900000}{\liter\per\text{døgn}}} \\
Befolkningsdekning ($P_{\text{dekket}}$) & Ved \textsc{dsb} basalrasjon (\SI{3}{\liter}/person/døgn) & \num{300000} pers. \\
Sikkerhetsfaktor ($SF$) & By med \num{75000} innbyggere & 4.0x \\
\bottomrule
\end{tabularx}
\end{table}

\subsection{To-spors operasjonsmodell (Triage-distribusjon)}
Flåten opererer etter en todelt fordelingsmodell:
\begin{itemize}[leftmargin=1.5em, itemsep=0.2em]
    \item \textbf{Linje A: Kritiske helseinstitusjoner (30 \% / 5 biler):} Prioritert leveranse til akuttsykehus og dialyseenheter. Bilene pumper vann direkte inn i byggenes reservebassenger på under 20 minutter via egne trykkpumper, uten manuelt bærearbeid.
    \item \textbf{Linje B: Desentraliserte sivile knutepunkter (70 \% / 10 biler):} Plasseres ved skolegårder og idrettshaller. Med AquaFeed-riggens 10 tappepunkter ekspederes opptil 800 personer i timen per bil. En rullerende stafett hindrer køkollaps.
\end{itemize}
